\documentclass[10pt,a4paper]{amsart}
\usepackage[margin=19mm]{geometry}
\usepackage{amsmath,amssymb,mathtools}
\usepackage[scr=rsfs,cal=euler]{mathalfa}
\usepackage{xfrac}
\usepackage[unicode,hidelinks,bookmarks=false]{hyperref}
\newcommand{\ie}{i.e.\ }
\newcommand{\cf}{cf.\ }
\newcommand{\resp}{respectively\ }
\newcommand{\Subsection}{\subsection}
\providecommand{\nobreakdash}{\nobreak\hbox{-}}
\setlength{\emergencystretch}{2em}
\multlinegap0pt
\usepackage{bm}
\usepackage{bbm}
\usepackage{dsfont}
\usepackage{xspace}
\usepackage{cancel}
\usepackage{mathtools}
\usepackage{amsthm}
\makeatletter
\@ifundefined{theorem}{\newtheorem{theorem}{Theorem}}{}
\@ifundefined{lemma}{\newtheorem{lemma}[theorem]{Lemma}}{}
\makeatother
\makeatletter
\newcommand{\eqnref}[1]{~{\textrm(\ref{#1})}}
\expandafter\ifx\csname proof\endcsname\relax
\renewenvironment{proof}{{\bfseries Proof.}}{\qed}
\fi
\makeatother
\title[Asymptotic growth of the number of Reciprocal Classes in the]{Asymptotic growth of the number of Reciprocal Classes in the Hecke Groups}
\author{Debattam Das, Krishnendu Gongopadhyay}
\date{Source: arXiv:2506.08449v1 (CC BY 4.0)}
\begin{document}
\maketitle

\noindent\emph{Source and licence.} Asymptotic growth of the number of Reciprocal Classes in the Hecke Groups. Version arXiv:2506.08449v1 (CC BY 4.0), distributed under the Creative Commons Attribution 4.0 International licence, as recorded in the arXiv licence metadata for this exact version and shown on the abstract page https://arxiv.org/abs/2506.08449v1. Full rights notice: licenses/arxiv-2506.08449-CC-BY-4.0.txt.\par\smallskip

\noindent\textit{Editorial scope.} Counted unit: the Lemma whose source label is \texttt{4.3} in the article (source0.tex lines 369--377, source label \texttt{4.3}), delivered together with the article's own complete original proof of it (source0.tex lines 378--429, one \texttt{proof} environment of 52 lines). No mathematics is written, repaired or re-derived: statement and proof are the article's own text line for line. Required context retained uncounted: none. Every definition, display and label of those lines is retained unchanged. Omitted from this package and not redistributed: 1--368, 430--660. Nothing inside a retained excerpt is omitted or altered: every retained body is delivered exactly as the article's lines are written, so no omission receipt of any kind accompanies this package. Citations: the retained excerpts carry no citation command at all, so no bibliography entry of the article is transcribed and no reference is redistributed. Reference adaptation: none was needed and none was made; every reference inside the delivered unit resolves inside it, so no unresolved reference or citation is produced. Printed slips kept literal: no printed word, symbol, hypothesis or number of the counted statement or of its proof was altered. Layout and class: the article's own class is \texttt{amsart} with packages including geometry, amsmath, amsfonts, amssymb, xcolor, url; the wrapper is the lane's pinned \texttt{amsart} editorial wrapper at 10pt on A4, which restates the theorem-like environments the retained text needs and adds the article's own macro definitions that the retained text needs (\texttt{\textbackslash eqnref}), with extra wrapper packages (bm, bbm, dsfont, xspace, cancel, mathtools). The delivered numbering is the unit's own. Assets: no retained span contains a figure environment, a graphics-inclusion command, a PDF-page inclusion, a table or a TikZ picture; no image file is copied into this package, the package declares no asset dependency and no image file is redistributed with it. Licence: version arXiv:2506.08449v1 of this article is distributed under the Creative Commons Attribution 4.0 International licence, as recorded in the arXiv licence metadata for this exact version and shown on the abstract page https://arxiv.org/abs/2506.08449v1. A full rights notice is delivered at licenses/arxiv-2506.08449-CC-BY-4.0.txt. Characters: every retained excerpt is pure ASCII, so the delivered engine, XeLaTeX with its default Latin Modern text font, reports no missing character.
	 \begin{lemma}\label{4.3}
 Consider the following equation: 	
	 	\begin{equation}
	 		\sum_{i=1}^{n} |~p_{i}~|=S_{n}, 
	 	\end{equation}
	 	where, $p_{i}$, and $x$ are positive integers, and  \( -r \leq p_i \leq r \). Let $\mathcal{N}(S_{n}\leq x) $ denote the number of solutions of the above equation subject to the condition $S_n \leq x$.  Then,  \[ \mathcal{N}(S_{n}\leq x)\sim \sum_{n=1}^{\lfloor \frac{x}{r}\rfloor}(2r)^{n}+{\sum_{n=\lceil \frac{x}{r}\rceil}^{x}(2r)^{n}\Phi\left(\frac{x-n\mu}{\sqrt{n}\sigma}\right)} .\]
	 	
	 	where, ${\mu=\frac{r+1}{2}}, ~\sigma^{2}=\frac{r^2}{12}$ and $\Phi$ is the cumulative distribution function of the standard Normal distribution $N(0,1).$
	 \end{lemma}

	 \begin{proof} 	
	 	Let $B:=\{-r,-r+1,\dots,-1,1,\dots,r\}$. 	Consider the inequality,
	 	\begin{equation} \label{eq:main2}
	 		\sum_{i=1}^{n} |p_i| \leq x,
	 	\end{equation}
	 	where each \( p_i \in \mathbb{Z} \setminus \{0\} \), and \( -r \leq p_i \leq r \), i.e., \( p_i \in B \). Suppose $\lceil\frac{x}{r}\rceil \leq n\leq x. $
	 	
	 	Let \( X_i = |p_i| \in \{1, 2, \dots, r\} \). From \eqnref{eq:main2}, 
	 	\[
	 	\sum_{i=1}^n X_i \leq x.
	 	\]
	 	
	 	Now consider \( X_i \) as independent and identically distributed random variables uniformly distributed on \(B \). This implies,  $\mathbb{P}(X_{i}=k)=\frac{2}{2r}=\frac{1}{r}$ for $k\in \{1,\dots,r\}$. Then
	 	\[
	 	\mathbb{E}[X_i] = \mu = \sum_{k=1}^{r} \frac{1}{r}\cdot k 
	 	\]
	 
	 	Therefore,	\[{\mu=\frac{r+1}{2}}.\]
	 	
	 	
	 	
	 	
	 	Also,
	 	
	 	\[ \qquad \mathrm{Var}(X_i) = \sigma^2 = \mathbb{E}[X^{2}]-\mu^{2}=\sum_{k=1}^{r} \frac{1}{r}\cdot k^{2} -\mu^{2}\]
	 	
	 	\[
	 	= \frac{1}{r} \cdot \frac{(r+1)r(2r + 1)}{6}  -\mu^{2}= \frac{(2r+ 1)(r+1)}{6} -\left( \frac{r+1}{2} \right)^2.
	 	\]
	 	

    It follows 	 	\[
	 	{
	 		\sigma^{2}=\frac{r^{2}-1}{12}.
	 	}
	 	\]
	 	
	 	
	 The total number of choices of the $n$-tuples  $(X_i)_{i=1}^{n}$ is $(2r)^{n}.$ Then
	 	by continuing the same process as the previous lemma, we obtain that, 
	 	 	the number of the solutions for $\sum_{i=1}^{n}|p_{i}|\leq x$ with $\lceil \frac{x}{r}\rceil\leq n\leq x$ is,
	 	\[ {\sum_{n=\lceil \frac{x}{r}\rceil}^{x}\mathcal{N}_{n}(S_n \leq ~x)}\sim{\sum_{n=\lceil \frac{x}{r}\rceil}^{x}(2r)^{n}\Phi\left(\frac{x-n\mu}{\sqrt{n}\sigma}\right)}. \]
	 	
	 	Now, if $n\leq \frac{x}{r}$ then any $n-$tuple of $X_{i}$'s will be part of the solutions of the Equation \eqnref{eq:main2}. Then,
	 	
	 	\[{\sum_{n=1}^{\lfloor \frac{x}{r}\rfloor}\mathcal{N}_{n}(S_n \leq ~x)}=\sum_{n=1}^{\lfloor \frac{x}{r}\rfloor}(2r)^{n}.\]
	 	
	 	Therefore, we have,
	 	The approximation of the total number of solutions,
	 	\[\sum_{n=1}^{\lfloor \frac{x}{r}\rfloor}(2r)^{n}+{\sum_{n=\lceil \frac{x}{r}\rceil}^{x}(2r)^{n}\Phi\left(\frac{x-n\mu}{\sqrt{n}\sigma}\right)} .\]
        This proves the lemma. 
	 \end{proof}

\end{document}
